\section{Nonlinear endpoint coordinates: a finite exact law}
\label{xid:sec:nonlinear}

The two recent dilation reports explicitly ask for the replacement of
their affine coordinate $pt$ by $pt+q_2t^2+q_3t^3+\cdots$
\cite{IncomingMellin,IncomingIdentities}.  This section answers that
question for every Stieltjes index and every argument derivative.  The
answer is local: a singularity at any other point is treated after
translation.  All finite parts use a right-hand cutoff in the indicated
coordinate.  Two-sided symmetric prescriptions are different objects.

Let $f$ be a real analytic local diffeomorphism with
\begin{equation}\label{xid:eq:f-germ}
 f(0)=0,\qquad f'(0)=p>0,\qquad
 H_f(t)=\frac{f(t)}t=p+q_2t+q_3t^2+\cdots.
\end{equation}
Choose a sufficiently small neighborhood so that $H_f$ is positive on
the real axis.  The power $H_f(t)^z$ uses the analytic logarithm whose
value at zero is $\log p$.  For a distribution $T$ the scalar pullback is
\begin{equation}\label{xid:eq:pullback-definition}
 \langle f^*T,\varphi\rangle
 =\left\langle T,
 \frac{\varphi(f^{-1}(y))}{f'(f^{-1}(y))}\right\rangle.
\end{equation}
Thus $f^*T=T\circ f$ for ordinary functions.  The Jacobian in this
definition is essential.

\subsection{A general coefficient lemma}

Write $\delta^{(j)}$ for the $j$th distributional derivative at zero,
so that $\langle\delta^{(j)},\varphi\rangle=(-1)^j\varphi^{(j)}(0)$.
For integers $q\ge1$ and $n\ge0$ set
\[
 U_{q,n}=\FP_y\bigl[\mathbf1_{y>0}y^{-q}\log^n y\bigr].
\]
Only the germ at zero matters; a smooth cutoff may be inserted away
from zero in every formula below.

\begin{theorem}[Coordinate defect for a logarithmic pole]
\label{xid:thm:coordinate-monomial}
One has
\begin{equation}\label{xid:eq:coordinate-monomial}
 f^*U_{q,n}
 =\FP_t\bigl[\mathbf1_{t>0}f(t)^{-q}\log^n f(t)\bigr]
       +\mathcal C_f(q,n),
\end{equation}
where the complete correction is the finite distribution
\begin{equation}\label{xid:eq:C-generator}
 \mathcal C_f(q,n)
 =n!\sum_{j=0}^{q-1}\frac{(-1)^j}{j!}
   [z^{n+1}t^{q-1-j}]H_f(t)^{z-q}\,\delta^{(j)}.
\end{equation}
It depends only on $f'(0),\ldots,f^{(q)}(0)$.
\end{theorem}

\begin{proof}
Begin with the locally integrable meromorphic family
$y_+^{z-q}$ for $\Re z>q-1$.  Pairing its pullback with a test function
gives
\[
 \int_0^h t^{z-q}H_f(t)^{z-q}\varphi(t)\,dt.
\]
Subtract the Taylor polynomial in $t$ through degree $q-1$ from the
holomorphic factor $H_f(t)^{z-q}\varphi(t)$.  The remainder is locally
normally integrable for $\Re z>-1$.  A Taylor term of degree $a$ has
integral
\[
 g_a(z)\frac{h^{z-q+a+1}}{z-q+a+1}.
\]
Near $z=0$, only $a=q-1$ has a pole.  The other Taylor terms have
nonzero denominators at zero, and taking their spectral coefficients
commutes with the coordinate finite part.  For the exceptional term,
the cutoff expression is
\[
 g_{q-1}(z)\frac{h^z-\varepsilon^z}{z}.
\]
The constant term in $\varepsilon$ of its coefficient of $z^n$ is the
regular coefficient of $g_{q-1}(z)h^z/z$, less
$[z^{n+1}]g_{q-1}(z)$.  In the original $y$ coordinate that last
correction is zero, because its Taylor coefficient is independent of
$z$.  Therefore the difference after pullback is
\[
 n![z^{n+1}t^{q-1}]H_f(t)^{z-q}\varphi(t).
\]
Expanding the test function and using the convention for
$\delta^{(j)}$ proves \eqref{xid:eq:C-generator}.  Its largest required
power of $t$ is $q-1$, proving the finite dependence assertion.
\end{proof}

Every function with a finite expansion in negative integer powers and
nonnegative logarithmic powers is covered by linearity.  Terms
$t^a\log^b t$ with $a\ge0$ are locally integrable and create no defect.
Consequently the theorem also applies to products of singular germs
after their ordinary local expansion has been performed; it does not
require a multiplication of distributions.

\begin{corollary}[A sharp disappearance threshold]
\label{xid:cor:monomial-threshold}
If $f'(0)=1$, then $\mathcal C_f(q,n)=0$ for $n\ge q-1$.
For $q\ge2$, writing $f(t)=t+a t^2+O(t^3)$, one has
\begin{equation}\label{xid:eq:monomial-top}
 \mathcal C_f(q,q-2)=\frac{a^{q-1}}{q-1}\,\delta.
\end{equation}
Thus the threshold is attained when $a\ne0$.
\end{corollary}

\begin{proof}
When $H_f(0)=1$, $[t^k]H_f(t)^{z-q}$ is a polynomial in $z$ of degree
at most $k$.  This follows by expanding
$\exp((z-q)\log H_f(t))$, since $\log H_f(t)=O(t)$.
In \eqref{xid:eq:C-generator}, $k\le q-1$, so its coefficient of
$z^{n+1}$ vanishes for $n\ge q-1$.  At $n=q-2$, only $j=0$ can
survive; the leading coefficient of $[t^{q-1}]H_f(t)^{z-q}$ is
$a^{q-1}/(q-1)!$.  Multiplication by $(q-2)!$ proves the formula.
\end{proof}

\subsection{Composition is an exact cocycle}

For any logarithmic singular germ $u$ define
\[
 C_f(u)=f^*(\FP u)-\FP(u\circ f).
\]
Theorem~\ref{xid:thm:coordinate-monomial} computes this distribution from the
finite singular expansion of $u$.  The expansion of $u\circ f$ is
again of the same kind, up to locally integrable terms.

\begin{theorem}[Composition and explicit delta pullback]
\label{xid:thm:cocycle}
For orientation-preserving analytic local coordinates $f,g$,
\begin{equation}\label{xid:eq:cocycle}
 C_{f\circ g}(u)=g^*C_f(u)+C_g(u\circ f).
\end{equation}
The delta derivatives in this formula can be computed without series
reversion:
\begin{equation}\label{xid:eq:delta-pullback}
 f^*\delta^{(r)}
 =\sum_{j=0}^r\frac{(-1)^{r-j}r!}{j!}
     [t^{r-j}]H_f(t)^{-r-1}\,\delta^{(j)}.
\end{equation}
\end{theorem}

\begin{proof}
Pullbacks satisfy $(f\circ g)^*=g^*f^*$.  Insert and subtract
$g^*\FP(u\circ f)$ in the definition of $C_{f\circ g}$; this gives
\eqref{xid:eq:cocycle} exactly.  For \eqref{xid:eq:delta-pullback}, take the
residue at $z=0$ in the proof of Theorem~\ref{xid:thm:coordinate-monomial}
with $q=r+1$.  The residue of $y_+^{z-r-1}$ is
$(-1)^r\delta^{(r)}/r!$.  The residue after pullback pairs with
$\varphi$ as $[t^r]H_f(t)^{-r-1}\varphi(t)$.
Expanding $\varphi$ gives the displayed formula.
\end{proof}

This proves the composition requirement in the source question, rather
than merely checking it for low-degree maps.  The accompanying exact
code also compares both sides using independently composed rational
Taylor series.

\subsection{Every Stieltjes index and every argument derivative}

Let $G_n$ be the canonical periodic finite part of $\gamma_n$ in the
unit coordinate.  Its local meromorphic generator is
\begin{equation}\label{xid:eq:periodic-G-local}
 Z(1-z)=\frac{\delta_0-1}{z}
          +\sum_{n\ge0}G_n\frac{z^n}{n!}.
\end{equation}
Indeed, locally across the endpoint,
$\zeta(1-z,\{y\})=y_+^{z-1}+\zeta(1-z,1+y)$.
The second term is smooth in $y$ and has only the constant spectral
pole $-1/z$.  Set
\begin{equation}\label{xid:eq:Er}
 E_r(z)=\prod_{h=1}^r(1-z/h),\qquad E_0(z)=1.
\end{equation}

\begin{theorem}[Nonlinear Stieltjes contact law]
\label{xid:thm:stieltjes-nonlinear}
For all $n,r\ge0$, as local distributions,
\begin{equation}\label{xid:eq:stieltjes-nonlinear}
 f^*(D_y^rG_n)
 =\FP_t\bigl[\gamma_n^{(r)}(\{f(t)\})\bigr]
       +A_{n,r}(f),
\end{equation}
where
\begin{equation}\label{xid:eq:A-generator}
 A_{n,r}(f)=n!r!\sum_{j=0}^r\frac{(-1)^{r+j}}{j!}
 [z^{n+1}t^{r-j}]
 E_r(z)H_f(t)^{z-r-1}\,\delta^{(j)}.
\end{equation}
The correction depends only on the $(r+1)$-jet of $f$.
\end{theorem}

\begin{proof}
The singular part of $D_y^r Z(1-z)$ is
$(-1)^r(1-z)_r y_+^{z-r-1}$.  After pullback, pair it with a test
function and repeat the Taylor subtraction proof above.  Only its
degree-$r$ Taylor coefficient creates a pole at $z=0$; the correction
is
\[
 n![z^{n+1}t^r](-1)^r(1-z)_r
                  H_f(t)^{z-r-1}\varphi(t).
\]
Since $(1-z)_r=r!E_r(z)$, this is \eqref{xid:eq:A-generator}.
The smooth local Hurwitz term commutes with coefficient extraction
and contributes no contact distribution.  The largest Taylor order
needed is $r$, which proves the jet assertion.
\end{proof}

For an affine germ $f(t)=pt$, only the term $j=r$ remains, and
\begin{equation}\label{xid:eq:A-affine}
 A_{n,r}(pt)=p^{-r-1}b_{n,r}(\log p)\delta^{(r)},\qquad
 b_{n,r}(L)=n![z^{n+1}]e^{Lz}E_r(z).
\end{equation}
After multiplication by $p^r$, this is precisely the earlier
ambient-derivative dilation law.  In particular,
\[
 b_{0,r}(L)=L-H_r,\qquad
 b_{1,r}(L)=\frac{L^2}{2}-H_rL+
                    \frac{H_r^2-H_r^{(2)}}2.
\]
The coefficient algebra is therefore an explicit extension of the
source's harmonic contact polynomial, with the same normalization.

For $f(t)=pt+q_2t^2+O(t^3)$ and $L=\log p$, the first two argument
orders take the particularly small form
\begin{align}
 A_{n,0}(f)&=\frac{L^{n+1}}{(n+1)p}\delta,\label{xid:eq:A0}\\
 A_{n,1}(f)&=\frac{n!}{p^2}[z^{n+1}](1-z)e^{Lz}\delta'
 +\frac{n!q_2}{p^3}[z^{n+1}](2-3z+z^2)e^{Lz}\delta.
 \label{xid:eq:A1}
\end{align}
Curvature creates an actual lower-order contact term; replacing $f$
by its tangent line loses it.

\begin{corollary}[Sharp unit-tangent Stieltjes threshold]
\label{xid:cor:stieltjes-threshold}
Suppose $f(t)=t+a t^2+O(t^3)$ and $r\ge1$.  Then
\begin{equation}\label{xid:eq:stieltjes-threshold}
 A_{n,r}(f)=0\quad(n\ge2r),\qquad
 A_{2r-1,r}(f)=\frac{(2r-1)!}{r!}a^r\delta.
\end{equation}
\end{corollary}

\begin{proof}
The degree in $z$ of
$[t^{r-j}]E_r(z)H_f(t)^{z-r-1}$ is at most $2r-j$.
This proves the vanishing assertion.  At $n=2r-1$ only $j=0$
survives.  The leading coefficients are $(-1)^r/r!$ for $E_r$
and $a^r/r!$ for $[t^r]H_f^{z-r-1}$.  Substitute them into
\eqref{xid:eq:A-generator} to obtain the second formula.
\end{proof}

For $f(t)=t+a t^2+b t^3+O(t^4)$ some exact coefficients are
\begin{align}
 A_{0,1}&=-\delta'-3a\delta,& A_{1,1}&=a\delta,\notag\\
 A_{0,2}&=-\tfrac32\delta''-11a\delta'
                       +(11b-25a^2)\delta,&
 A_{3,2}&=3a^2\delta.\label{xid:eq:A-table}
\end{align}
Together with \eqref{xid:eq:stieltjes-threshold}, these examples show why
the earlier zero contact terms at unit affine scale cannot simply be
transported through a nonlinear coordinate.

\subsection{Exact antiderivatives and their endpoint constants}

There is a useful scalar consequence which does not require any
distribution notation.  Suppose $f$ is increasing on $[0,b]$,
$f(0)=0$, $f'(0)>0$, and $B=f(b)>0$.  Assume that the Taylor germ at
zero is as above.  For $r\ge1$ define
\begin{equation}\label{xid:eq:primitive-E}
 \mathcal E_{n,r}(f)=\gamma_n^{(r-1)}(1)
 +(-1)^{r-1}n![z^nt^r](1-z)_{r-1}H_f(t)^{z-r}.
\end{equation}

\begin{theorem}[Finite-part primitive endpoint identity]
\label{xid:thm:nonlinear-primitive}
For every $n\ge0$ and $r\ge1$,
\begin{equation}\label{xid:eq:nonlinear-primitive}
 \FP_t\int_0^b f'(t)\gamma_n^{(r)}(f(t))\,dt
       =\gamma_n^{(r-1)}(B)-\mathcal E_{n,r}(f).
\end{equation}
For argument order zero the exact change-of-coordinate identity is
\begin{equation}\label{xid:eq:primitive-order0}
 \FP_t\int_0^b f'(t)\gamma_n(f(t))\,dt
 =\FP_y\int_0^B\gamma_n(y)\,dy
       -\frac{\log^{n+1}p}{n+1}.
\end{equation}
\end{theorem}

\begin{proof}
At a positive cutoff, the fundamental theorem of calculus gives
$\gamma_n^{(r-1)}(B)-\gamma_n^{(r-1)}(f(\varepsilon))$.
Use
\[
 \gamma_n(y)=y^{-1}\log^n y+\gamma_n(1+y)
\]
and differentiate $r-1$ times.  The smooth term has constant
$\gamma_n^{(r-1)}(1)$.  Write its singular term as the coefficient
of $z^n/n!$ in
$(-1)^{r-1}(1-z)_{r-1}t^{z-r}H_f(t)^{z-r}$.
The constant term in $t$ comes only from the Taylor term of degree
$r$, with no power of $\log t$ retained.  This is exactly
\eqref{xid:eq:primitive-E}.  For $r=0$, substitution gives a lower
endpoint primitive $\log^{n+1}f(\varepsilon)/(n+1)$, whose constant
term is $\log^{n+1}p/(n+1)$.  All smooth endpoint contributions agree
in the two coordinates, proving \eqref{xid:eq:primitive-order0}.
\end{proof}

For example, put $f(t)=t+c t^2$ on an interval where it is increasing.
Then $B=b+cb^2$ and the formulas give
\begin{align}
 \FP\int_0^b(1+2ct)\psi'(t+ct^2)\,dt
    &=\psi(B)+\gamma-c,\label{xid:eq:psi-prime-curvature}\\
 \FP\int_0^b(1+2ct)\psi''(t+ct^2)\,dt
    &=\psi'(B)-\zeta(2)-3c^2,\label{xid:eq:psi-double-curvature}\\
 \FP\int_0^b(1+2ct)\gamma_1'(t+ct^2)\,dt
    &=\gamma_1(B)-\gamma_1(1)-c.\label{xid:eq:gamma1-curvature}
\end{align}
These are exact antiderivative identities with explicit curvature
constants.  In particular, a formal substitution which silently
equates the $t$ and $f(t)$ finite parts gives an incorrect constant.

The zero-order Hurwitz primitive on the right of
\eqref{xid:eq:primitive-order0} can be written, if desired, by integrating
$\zeta(1-z,y)$ as $\zeta(-z,y)/z$ and taking its regular spectral
coefficients.  At $n=0$, $\gamma_0=-\psi$ gives the familiar
logarithmic Gamma primitive.  The new ingredient above is the exact
endpoint transformation, not the classical Hurwitz primitive itself
\cite{EspinosaMoll2}.
